% =============================================================
% บทที่ 5 - สรุปผลการดำเนินงาน
% โครงงานประเภท B (งานวิจัย) — สรุปตามวัตถุประสงค์
% เทียบกับเส้นฐาน ตอบคำถามวิจัย
% =============================================================

% -----------------------------------------------------------
\section{สรุปผลการดำเนินโครงงาน}
% -----------------------------------------------------------

\subsection{ผลตามวัตถุประสงค์แต่ละข้อ}

สรุปผลเทียบกับวัตถุประสงค์ที่ตั้งไว้ในบทที่ 1 แสดงดังตารางที่ \ref{tab:obj_summary}

\begin{table}[H]
\centering
\caption{สรุปผลตามวัตถุประสงค์}
\label{tab:obj_summary}
\begin{tabular}{|c|p{5.4cm}|c|p{5.0cm}|}
\hline
\textbf{ข้อ} & \textbf{วัตถุประสงค์} & \textbf{สถานะ} & \textbf{หลักฐาน} \\
\hline
1 & พัฒนาระบบแนะนำสินค้าสองภาษาบนเครื่องส่วนบุคคล & บรรลุ & รันครบ 840 คำตอบ ไม่มีข้อใดทำงานผิดพลาด \\ \hline
2 & เปรียบเทียบ RAG กับ Agentic RAG แบบเปลี่ยนตัวแปรทีละตัว & บรรลุ & ทุกสถาปัตยกรรมใช้ส่วนประกอบร่วมกัน ต่างเฉพาะชั้นตัดสินใจ \\ \hline
3 & ศึกษาว่าความสามารถตัดสินใจเองช่วยในสถานการณ์ใด & บรรลุ & ช่วยเฉพาะกลุ่มที่คลังไม่มีสินค้าตรง ไม่ช่วยเลยเมื่อคลังมีของ \\ \hline
4 & สร้างและเผยแพร่ชุดทดสอบพร้อมเฉลย & บรรลุ & เผยแพร่คำค้น 140 คำ เฉลยความตรง 591 คู่ และเฉลยความซื่อตรง 402 รายการ \\
\hline
\end{tabular}
\end{table}

\subsection{คำตอบของคำถามวิจัย}

\textbf{คำถามที่ 1 สถาปัตยกรรมแบบ Agentic RAG ให้ผลดีกว่าแบบ RAG มาตรฐานหรือไม่}

ตอบว่า\textbf{ดีกว่าทั้งสองภาษา และไม่แพ้แม้แต่ข้อเดียว} ฝั่งอังกฤษ nDCG@3 เพิ่มจาก 0.3288 เป็น 0.3564 ดีขึ้น 8 ข้อ แย่ลง 0 ข้อ ที่ $p = 0.0078$ ฝั่งไทยเพิ่มจาก 0.3165 เป็น 0.3267 ดีขึ้น 5 ข้อ แย่ลง 0 ข้อ ที่ $p = 0.0625$ ซึ่งเป็นค่าต่ำสุดที่เป็นไปได้เมื่อมีข้อที่ต่างกัน 5 ข้อ

\textbf{คำถามที่ 2 ความได้เปรียบนั้นมาจากไหน}

ตอบว่า\textbf{มาจากกลุ่มที่คลังไม่มีสินค้าตรงเพียงกลุ่มเดียว} ในกลุ่มที่คลังมีสินค้าตรง ทั้งสามสถาปัตยกรรมได้คะแนนเท่ากันทุกทศนิยมในฝั่งอังกฤษ และสถิติการทำงานยืนยันว่าการค้นซ้ำไม่ทำงานเลยเมื่อคลังมีของ ข้อสรุปเชิงปฏิบัติคือ\textbf{ความเป็นเอเจนต์คุ้มค่าเมื่อคำค้นมีโอกาสหาของไม่เจอ ไม่ใช่ทุกกรณี}

\textbf{คำถามที่ 3 ควรขยายขอบเขตการตัดสินใจของเอเจนต์ไปถึงการคุมการค้นซ้ำหรือไม่}

ตอบว่า\textbf{ไม่คุ้ม} เพราะ A3C กับ A3 ให้ผล nDCG เท่ากันทุกทศนิยมทั้งสองภาษา แต่ A3 ใช้การเรียกแบบจำลองมากกว่าประมาณ 1.3 ครั้งต่อข้อ และใช้เวลามากกว่าประมาณ 5 วินาทีต่อข้อ คู่นี้เป็นคู่เดียวที่เปลี่ยนตัวแปรทีละตัวอย่างแท้จริง คำตอบจึงเชื่อถือได้

\subsection{ข้อค้นพบที่สำคัญที่สุด}

ข้อค้นพบที่มีผลต่อการออกแบบระบบมากที่สุด\textbf{ไม่ใช่เรื่องสถาปัตยกรรม แต่เป็นเรื่องชั้นจัดอันดับ} การตรวจสอบทุกขั้นของระบบพบว่าชั้นจัดอันดับทำสินค้าที่ดีสูญหายไป 45 ถึง 48 เปอร์เซ็นต์ ซึ่งมากกว่าทุกชั้นอื่นรวมกันประมาณ 3.5 เท่า

เมื่อเพิ่มสัญญาณการนับคำที่ซ้อนกับคำค้นเข้าไปในชั้นจัดอันดับ ได้ผลเพิ่มขึ้น 3.4 ถึง 3.5 จุดแบบตลอดสาย ซึ่ง\textbf{มากกว่าผลที่ความเป็นเอเจนต์ให้ประมาณ 2 เท่า} ข้อค้นพบนี้ชี้ว่าในระบบลักษณะนี้ การปรับปรุงส่วนประกอบพื้นฐานอาจให้ผลตอบแทนสูงกว่าการเพิ่มความซับซ้อนของสถาปัตยกรรม

% -----------------------------------------------------------
\section{ข้อจำกัดของโครงงาน}
% -----------------------------------------------------------

\begin{enumerate}[label=\arabic*)]
    \item \textbf{ขนาดของแบบจำลองภาษา} ใช้แบบจำลองขนาด 8 พันล้านพารามิเตอร์ ซึ่งเป็นขนาดใหญ่ที่สุดที่ทำงานได้บนหน่วยความจำกราฟิก 8 กิกะไบต์ ข้อสรุปที่ว่าแบบจำลองตัดสินความตรงไม่ได้จึงผูกกับขนาดนี้ และสรุปไปถึงแบบจำลองขนาดใหญ่กว่าไม่ได้
    \item \textbf{ไม่มีผู้ติดป้ายอิสระคนที่สอง} เฉลยทั้งหมดติดโดยผู้วิจัยซึ่งเป็นผู้พัฒนาระบบที่ถูกประเมินเอง จึงยังรายงานค่าความสอดคล้องระหว่างผู้ติดป้ายไม่ได้ และยืนยันไม่ได้ว่าเกณฑ์ที่ใช้ตรงกับที่ผู้อื่นจะตัดสิน
    \item \textbf{ความเอนเอียงของกองที่ติดเฉลย} กองที่ติดเฉลยมาจากสินค้าที่ระบบของโครงงานนี้หยิบขึ้นมาเท่านั้น ระบบอื่นที่หยิบสินค้านอกกองจะถูกประเมินต่ำกว่าความเป็นจริง
    \item \textbf{ตัวหารของ nDCG ไม่เป็นมาตรฐาน} ใช้ระบบในอุดมคติเป็นตัวหารแทนกองที่ติดเฉลย ค่าที่ได้จึงเทียบกับงานวิจัยอื่นโดยตรงไม่ได้
    \item \textbf{คำค้นภาษาไทยเป็นคำแปล} ไม่ใช่คำที่ผู้ใช้ชาวไทยพิมพ์เอง ความยากจากการสะกดผิดจึงหายไปบางส่วน และในกลุ่มที่คลังไม่มีของ มีคำที่แปลจริงเพียง 15 จาก 60 คำ เพราะที่เหลือเป็นชื่อเฉพาะที่แปลไม่ได้ ทำให้กองไทยกลุ่มนี้ซ้อนทับกองอังกฤษถึง 75 เปอร์เซ็นต์ ตัวเลขฝั่งไทยในกลุ่มนี้จึงไม่ใช่การวัดภาษาไทยที่เป็นอิสระ
    \item \textbf{ชุดทดสอบไม่มีสถานะกลาง} ไม่ครอบคลุมกรณีที่คลังมีสินค้าใกล้เคียงแต่ไม่ตรง ซึ่งเป็นสถานการณ์ที่การค้นซ้ำน่าจะได้เปรียบมากที่สุด การประเมินจึงอาจต่ำกว่าประโยชน์จริง
    \item \textbf{กฎตัดคำของสัญญาณนับคำซ้อนกำหนดด้วยมือ} รายการคำที่ไม่ช่วยแยก 13 คำ และเกณฑ์ตัดคำที่สั้นกว่า 3 ตัวอักษร เขียนขึ้นเองโดยไม่ได้เรียนรู้จากข้อมูลและไม่ได้กวาดค่าเปรียบเทียบกับทางเลือกอื่น จึงยืนยันไม่ได้ว่าเป็นค่าที่ดีที่สุด และผูกกับโดเมนสินค้ากับภาษาอังกฤษ ผลข้างเคียงที่วัดได้คือคำค้นที่เขียนตัวเลขแยกจากหน่วยจะถูกตัดทิ้งทั้งคู่ ทำให้สัญญาณแยกความจุที่ต่างกันไม่ได้
    \item \textbf{ปรับค่าน้ำหนักบนชุดเดียวกับที่รายงานผล} น้ำหนักของสัญญาณการนับคำซ้อนเลือกจากการกวาดค่าบนเฉลยชุดเดียวกับที่ใช้รายงาน ไม่ได้แบ่งกองพัฒนาออกจากกองทดสอบในขั้นนี้
    \item \textbf{เฉลยความซื่อตรงเป็นของรุ่นก่อน} เฉลยความซื่อตรงของคำตอบวัดบนคำตอบที่ได้ก่อนเปลี่ยนชั้นจัดอันดับ เมื่อสินค้าที่หยิบเปลี่ยน คำตอบย่อมเปลี่ยนตาม ค่าที่รายงานจึงใช้กับระบบรุ่นปัจจุบันโดยตรงไม่ได้
\end{enumerate}

% -----------------------------------------------------------
\section{ปัญหา อุปสรรค และแนวทางแก้ไข}
% -----------------------------------------------------------

\subsection{การเปรียบเทียบที่เปลี่ยนหลายตัวแปรพร้อมกัน}

\textbf{ปัญหา} รุ่นแรกของการทดลองเปรียบเทียบ RAG กับ Agentic RAG โดยที่สองสถาปัตยกรรมต่างกันทั้งชั้นวางแผนและชั้นค้นซ้ำพร้อมกัน จึงสรุปไม่ได้ว่าผลมาจากส่วนใด

\textbf{แนวทางแก้ไข} เพิ่มสถาปัตยกรรมที่สามที่ต่างจากแบบเต็มเพียงตัวแปรเดียว ทำให้ได้คู่เปรียบเทียบที่เปลี่ยนตัวแปรทีละตัวอย่างแท้จริง

\subsection{กองที่ติดเฉลยใช้วัดระบบที่เปลี่ยนไปแล้วไม่ได้}

\textbf{ปัญหา} หลังเปลี่ยนชั้นจัดอันดับ สินค้าที่ระบบหยิบขึ้นมาเปลี่ยนไปถึง 51 จาก 80 คำค้น สินค้าใหม่ไม่มีเฉลย และวิธีคิดตัวชี้วัดตัดคำค้นที่เฉลยไม่ครบทิ้งทั้งข้อ เหลือคำค้นที่เทียบได้เพียง 21 จาก 50 คำ และเกือบทั้งหมดเป็นคำที่ไม่มีอะไรเปลี่ยน จึงวัดอะไรไม่ได้

\textbf{แนวทางแก้ไข} ติดเฉลยเพิ่มสำหรับสินค้าที่โผล่ขึ้นมาใหม่ จนชุดคำค้นที่เทียบได้ครบกลับมาเป็น 140 จาก 140 คำ และบันทึกข้อจำกัดนี้ไว้ในรายงาน

\subsection{การทดสอบที่ผ่านแต่ไม่ครอบคลุมเส้นทางจริง}

\textbf{ปัญหา} การเพิ่มสัญญาณใหม่เข้าไปในชั้นจัดอันดับทำให้การเรียกใช้งานบางเส้นทางขาดพารามิเตอร์ แต่การทดสอบเบื้องต้นผ่านเพราะกรณีทดสอบที่ใช้ไม่ได้วิ่งผ่านเส้นทางนั้น ปัญหาจึงไปปรากฏตอนรันจริงและทำให้การรันล้มเหลวจำนวนมาก

\textbf{แนวทางแก้ไข} กำหนดเป็นแนวปฏิบัติว่าการทดสอบต้องใช้กรณีที่\textbf{ทำให้เส้นทางใหม่ทำงานจริง} ไม่ใช่เพียงกรณีที่ไม่ทำให้โปรแกรมหยุด

\subsection{ความเข้าใจผิดเกี่ยวกับความสามารถของการค้นซ้ำ}

\textbf{ปัญหา} ตอนแรกเข้าใจว่าการค้นซ้ำช่วยเพราะแก้คำสะกดผิดได้ การทดสอบกับคำที่สะกดผิดพบว่ากู้คืนได้เพียง 3 จาก 12 คำ และกรณีที่สำเร็จเป็นชื่อแบรนด์ดังที่แบบจำลองรู้จักอยู่แล้ว

\textbf{แนวทางแก้ไข} ตรวจสอบสมมติฐานด้วยการวัดก่อนนำไปเขียนเป็นข้อสรุป และแก้คำอธิบายให้ตรงกับหลักฐาน คือการค้นซ้ำช่วยด้วยการเปลี่ยนไปใช้คำที่คลังสินค้าใช้เรียก

% -----------------------------------------------------------
\section{ข้อเสนอแนะในการพัฒนาต่อไป}
% -----------------------------------------------------------

\begin{enumerate}[label=\arabic*)]
    \item \textbf{หาผู้ติดป้ายอิสระคนที่สอง} มาติดเฉลยซ้ำบนชุดตรวจ 48 คู่ที่เตรียมไว้แล้ว เพื่อยืนยันว่าเกณฑ์ที่ใช้ตรงกับที่ผู้อื่นจะตัดสิน ซึ่งจะทำให้รายงานค่าความสอดคล้องระหว่างผู้ติดป้ายได้
    \item \textbf{เพิ่มกลุ่มคำค้นที่คลังมีสินค้าใกล้เคียงแต่ไม่ตรง} เพื่อวัดประโยชน์ของการค้นซ้ำในสถานการณ์ที่ควรได้เปรียบที่สุด ซึ่งชุดทดสอบปัจจุบันยังไม่ครอบคลุม
    \item \textbf{ติดเฉลยความซื่อตรงของคำตอบใหม่} บนคำตอบที่ได้จากระบบรุ่นปัจจุบัน เพื่อให้ตัวชี้วัดด้านนี้ตรงกับระบบที่รายงาน
    \item \textbf{ถ่วงน้ำหนักคำตามความถี่ผกผันในคลัง} แทนการใช้รายการคำที่ไม่ช่วยแยกแบบเขียนเอง เพื่อให้ระบบหาได้เองว่าคำใดไม่ช่วยแยกความตรง ซึ่งจะปรับตามโดเมนได้โดยอัตโนมัติ และแก้ปัญหาที่ตัวเลขกับหน่วยสั้นถูกตัดทิ้ง
    \item \textbf{ใช้คำค้นภาษาไทยที่ผู้ใช้ไทยพิมพ์เอง} แทนคำแปล เพื่อให้ความยากที่เกิดจากการสะกดและการเลือกคำสะท้อนการใช้งานจริง
    \item \textbf{ทดสอบกับแบบจำลองภาษาขนาดใหญ่กว่า} เพื่อตรวจว่าข้อสรุปที่ว่าแบบจำลองตัดสินความตรงไม่ได้นั้นเป็นข้อจำกัดของขนาด หรือเป็นข้อจำกัดของแนวทาง
    \item \textbf{ปรับปรุงชั้นแปลภาษา} เนื่องจากความคลาดเคลื่อนของชั้นนี้ไหลลงไปทำลายสัญญาณปลายน้ำ โดยเฉพาะสัญญาณการนับคำซ้อนซึ่งฝั่งไทยไม่ได้ประโยชน์เลย
    \item \textbf{พัฒนาส่วนติดต่อผู้ใช้} เพื่อให้ทดสอบกับผู้ใช้จริงได้ ซึ่งจะเปิดทางไปสู่การวัดความพึงพอใจและความมีประโยชน์ที่ตัวชี้วัดปัจจุบันยังวัดไม่ได้
\end{enumerate}
